\documentclass[12pt,a4paper]{report}
\usepackage[utf8]{inputenc}
\usepackage[T1]{fontenc}
\usepackage{geometry}
\usepackage{setspace}
\usepackage{hyperref}
\usepackage{enumitem}
\usepackage{xcolor}
\geometry{margin=1in}
\onehalfspacing
\setlist[itemize]{leftmargin=*,itemsep=0.25em}
\setlist[enumerate]{leftmargin=*,itemsep=0.25em}
\hypersetup{colorlinks=true,linkcolor=blue,urlcolor=blue,citecolor=blue}

\title{AI-Powered Unified Fashion Discovery Platform for the Lebanese Market}
\author{Student Name (ID): [Fill] \\ Supervisor Name: [Fill]}
\date{[Fill]}

\begin{document}
\maketitle
\tableofcontents
\clearpage


\chapter{Signature Page}
Student Signature: \_\_\_\_\_\_\_\_\_\_\_\_

Supervisor Signature: \_\_\_\_\_\_\_\_\_\_\_\_


\chapter{Acknowledgments}
The author acknowledges the guidance of the supervisor, the support of faculty members, and the collaborative contributions of teammates involved in backend engineering, frontend development, and system testing.


\chapter{Abstract}
This project presents the design and implementation of an AI-powered fashion aggregation platform that unifies fragmented product catalogs and enables inspiration-driven discovery. Traditional e-commerce systems are optimized for keyword queries and are weak at visual and outfit-based intent. To address this limitation, the platform combines multimodal retrieval (text and image), shop-the-look detection, recommendation and outfit completion, wardrobe intelligence, virtual try-on, and intelligent product comparison in one architecture.


The system is implemented as a modular Node.js/TypeScript API using Express, PostgreSQL, OpenSearch, Redis/BullMQ, and cloud/object storage integrations. The machine-learning stack integrates CLIP embeddings, YOLO-based detection, BLIP captioning, Gemini-assisted intent/attribute processing, XGBoost ranking, and Vertex AI virtual try-on. Results indicate strong practical capability for real-time product discovery and decision support, while also revealing known gaps (e.g., checkout workflow and certain auth hardening features). The project demonstrates a production-oriented, extensible architecture for local-market fashion intelligence.


\chapter*{References}
\addcontentsline{toc}{chapter}{References}

\chapter*{Appendices}
\addcontentsline{toc}{chapter}{Appendices}


\chapter{1. Introduction}
\section{1.1 Background and Motivation}
Fashion search behavior is shifting from text-only lookup toward visual and contextual discovery. Users increasingly discover outfits from social media and real-world observations, then seek similar purchasable items. In fragmented markets, this process is inefficient because products are distributed across separate stores and inconsistent metadata structures.


\section{1.2 Problem Statement}
The core problem is the absence of a unified, intelligent platform that can convert visual inspiration into actionable shopping options across multiple vendors. This includes:

\begin{itemize}
\item Fragmented catalogs
\item Limited visual search support
\item Weak cross-product decision support
\item Insufficient outfit-level reasoning
\item Performance constraints in real-time retrieval
\end{itemize}

\section{1.3 Objectives}
The project objectives are to:

\begin{itemize}
\item Aggregate fashion products into a unified index
\item Support multimodal search (text, image, and multi-image prompts)
\item Enable shop-the-look decomposition and retrieval
\item Provide intelligent comparison and recommendation features
\item Offer outfit completion and wardrobe-assisted styling workflows
\item Maintain scalable and low-latency system behavior
\end{itemize}

\section{1.4 Significance and Impact}
The platform reduces user effort, improves discovery relevance, increases merchant visibility, and demonstrates an applied software + AI solution tailored to a local ecosystem with strong extensibility toward regional scaling.


\chapter{2. Literature Review / Related Work}
\section{2.1 Traditional E-Commerce Search}
Conventional systems rely on keyword and metadata matching, which is effective for explicit queries but weak for visual intent and implicit style attributes.


\section{2.2 Multimodal Retrieval Systems}
Modern systems align image and text in shared embedding spaces to support similarity search. Practical quality depends on embedding fidelity, indexing strategy, and reranking logic.


\section{2.3 Fashion-Specific AI Systems}
Fashion introduces complex attributes (color harmony, texture, pattern, layering, style compatibility). Specialized pipelines often include detection, captioning, and recommendation signals.


\section{2.4 Gaps in Existing Approaches}
Most solutions remain fragmented: they focus on one capability (search, recommendation, or comparison) rather than an integrated end-to-end decision pipeline.


\section{2.5 Identified Gap and Proposed Solution}
This project addresses an integration gap by combining aggregation, multimodal discovery, intelligent compare, outfit completion, wardrobe intelligence, and virtual try-on within one production-oriented architecture.


\chapter{3. System Analysis and Design}
\section{3.1 Requirements Analysis}
\subsection{3.1.1 Functional Requirements}
\begin{itemize}
\item User authentication, profile, cart, and favorites
\item Text/image/multi-image/multi-vector search
\item Shop-the-look image analysis and per-item retrieval
\item Product recommendation, compare, and complete-style
\item Wardrobe management and style intelligence
\item Async virtual try-on workflow
\item Ingestion and labeling pipelines for data enrichment
\end{itemize}

\subsection{3.1.2 Non-Functional Requirements}
\begin{itemize}
\item Low response latency and throughput resilience
\item Scalable modular architecture
\item Fault tolerance with graceful degradation for optional ML services
\item Operational observability (health and metrics)
\item Security controls for protected routes and admin actions
\end{itemize}

\section{3.2 System Architecture}
\subsection{3.2.1 Architectural Style and Rationale}
The platform uses a modular service-layer monolith as the default runtime shape, with optional split-role deployment for operational scaling. This means one codebase and one API contract, while allowing different runtime responsibilities (\texttt{api}, \texttt{ml}, or \texttt{all}) based on environment configuration.

The selected style balances three goals:
\begin{itemize}
\item fast feature iteration in a capstone/product-discovery phase
\item strict module boundaries to reduce coupling and simplify testing
\item progressive path toward service decomposition without rewriting business contracts
\end{itemize}

The primary internal boundary is \texttt{routes -> controller -> service}, where:
\begin{itemize}
\item routes map endpoint paths and middleware
\item controllers normalize/validate request context and shape responses
\item services contain business logic, DB/search I/O, and orchestration logic
\end{itemize}

\subsection{3.2.2 Logical View (Layered Architecture)}
The logical architecture is organized into four layers:
\begin{itemize}
\item Presentation layer: storefront and admin clients consuming REST endpoints
\item API layer: Express middleware stack and route modules
\item Intelligence layer: retrieval, reranking, detection, recommendation, compare, and try-on orchestration
\item Data layer: PostgreSQL (transactional/system-of-record), OpenSearch (retrieval index), Redis/BullMQ (asynchronous jobs), object storage (images/artifacts)
\end{itemize}

The key architectural principle is separation of transaction-heavy workflows (auth/cart/favorites/wardrobe state) from compute-heavy workflows (search embeddings, detection, try-on generation), even when they share the same deployment unit.

\subsection{3.2.3 Runtime View (Request Lifecycle)}
A typical synchronous request follows:
\begin{enumerate}
\item middleware entry (security headers, CORS, parsing, logging, metrics, rate limiting)
\item controller-level validation and request normalization
\item service execution across PostgreSQL/OpenSearch/external ML clients as needed
\item response shaping to a stable API contract
\item centralized error middleware for normalized failure payloads
\end{enumerate}

For asynchronous workflows (ingest and try-on), the runtime pattern is:
\begin{enumerate}
\item API accepts request and persists job state
\item job is queued in Redis/BullMQ
\item worker performs heavy processing and stores outputs
\item client polls status endpoint until completion/failure
\end{enumerate}

This split prevents long-running ML tasks from blocking request threads and keeps latency predictable for user-facing endpoints.

\subsection{3.2.4 Module and Route Topology}
The route topology is domain-oriented rather than technical-function oriented. Major mounted domains include:
\begin{itemize}
\item \texttt{/search} and \texttt{/products} for discovery and catalog experiences
\item \texttt{/api/images} for detection and shop-the-look operations
\item \texttt{/api/wardrobe} and \texttt{/api/tryon} for personal styling workflows
\item \texttt{/api/compare} for decision intelligence
\item \texttt{/api/auth}, \texttt{/api/cart}, \texttt{/api/favorites} for account and commerce primitives
\item \texttt{/admin}, \texttt{/metrics}, and \texttt{/health} for operations and governance
\end{itemize}

Each module is responsible for its own controller/service logic while reusing shared core libraries (database pool, OpenSearch client, shared ranking/feature utilities, image/model adapters).

\subsection{3.2.5 Data Architecture and Responsibilities}
Data responsibilities are intentionally split by access pattern:
\begin{itemize}
\item PostgreSQL: authoritative relational state (users, product metadata references, wardrobe entities, try-on jobs, workflow tracking)
\item OpenSearch: denormalized, vectorized retrieval index optimized for BM25 + kNN hybrid search
\item Redis/BullMQ: ephemeral queueing and workflow coordination for asynchronous pipelines
\item Object storage: binary assets (product images, user uploads, generated outputs)
\end{itemize}

The architecture follows a ``source-of-truth + projection'' model: relational state remains authoritative, while OpenSearch acts as a query-optimized projection for low-latency retrieval.

\subsection{3.2.6 Intelligence Architecture (ML/AI Integration)}
The intelligence layer is orchestration-centric rather than model-centric. The API coordinates multiple model/services:
\begin{itemize}
\item embedding and similarity services for image/text retrieval
\item detection/captioning services for item decomposition and semantic enrichment
\item ranker service for recommendation and reranking
\item cloud inference for virtual try-on generation
\end{itemize}

A core architectural decision is graceful degradation:
\begin{itemize}
\item if optional signals fail (e.g., captioning/ranker), fallback heuristics keep endpoints operational
\item failures are isolated to signal quality when possible, not elevated to full endpoint outages
\end{itemize}

\subsection{3.2.7 Deployment Architecture}
The system supports two deployment patterns:
\begin{itemize}
\item Monolith deployment (\texttt{SERVICE\_ROLE=all}): simplest operational model for integrated environments
\item Split-role deployment (\texttt{SERVICE\_ROLE=api} and \texttt{SERVICE\_ROLE=ml}): separates request-facing endpoints from compute-heavy paths while preserving API contracts
\end{itemize}

This gives a staged scalability strategy: start with lower operational complexity, then isolate pressure points (ML-heavy traffic) when load or cost profile requires it.

\subsection{3.2.8 Security and Reliability by Design}
Architectural controls are embedded at the platform boundary:
\begin{itemize}
\item security middleware and rate limiting at ingress
\item auth/admin middleware for protected domain enforcement
\item health/readiness endpoints for orchestrator-level lifecycle decisions
\item metrics endpoint for runtime observability and SLO tracking
\item centralized error handling for consistent failure semantics
\end{itemize}

\subsection{3.2.9 Architectural Trade-offs}
Key trade-offs in this design include:
\begin{itemize}
\item Monolith simplicity vs microservice isolation: chosen approach favors faster integration and lower cognitive overhead during rapid feature expansion.
\item Shared runtime dependencies vs strict service autonomy: shared code paths accelerate delivery but require stronger operational guardrails.
\item Hybrid retrieval quality vs infrastructure complexity: BM25 + vector + reranking improves relevance but increases calibration and monitoring requirements.
\end{itemize}

Overall, the architecture is intentionally evolutionary: it is stable enough for production-like operation, while still adaptable for future decomposition into independently scaled services.

\begin{figure}[htbp]
\centering
\fbox{
\begin{minipage}{0.95\textwidth}
\small
\begin{verbatim}
                    +----------------------------------+
                    |     Presentation Layer           |
                    |  Marketplace UI / Admin UI       |
                    +----------------+-----------------+
                                     |
                                     v
                    +----------------------------------+
                    |         API Layer (Express)      |
                    |  Routes -> Controllers -> Services|
                    |  Auth, Search, Compare, Wardrobe |
                    +----------------+-----------------+
                                     |
                  +------------------+------------------+
                  |                                     |
                  v                                     v
     +-----------------------------+       +-----------------------------+
     |      Intelligence Layer     |       |        Data Layer           |
     |  CLIP / YOLO / BLIP /       |<----->| PostgreSQL (source of truth)|
     |  Gemini / XGBoost / Try-On  |       | OpenSearch (retrieval index)|
     |  Orchestration + Reranking  |       | Redis/BullMQ (async jobs)   |
     +-----------------------------+       | Object Storage (images/files)|
                  |                         +-----------------------------+
                  v
     +-----------------------------+
     | External ML/Cloud Services  |
     | Vertex AI / Model Endpoints |
     +-----------------------------+
\end{verbatim}
\end{minipage}
}
\caption{High-level system architecture showing layered responsibilities and major runtime interactions.}
\label{fig:system-architecture}
\end{figure}

\noindent Figure~\ref{fig:system-architecture} summarizes the platform decomposition and interaction boundaries used throughout implementation and deployment.


\section{3.3 Feasibility and Risk Analysis}
Technical feasibility is supported by mature infrastructure and tooling. Key architectural risks include external model dependency, data quality variability from aggregated sources, and configuration/secret management complexity in distributed deployment.


\chapter{4. Implementation}
\section{4.1 Development Methodology}
Implementation follows iterative delivery with modular boundaries (\texttt{routes} -> \texttt{controller} -> \texttt{service}) and continuous integration of search, recommendation, wardrobe, compare, and try-on capabilities. The architecture prioritizes testable service layers, explicit API contracts, and feature-level observability.


\section{4.2 Tools, Technologies, and Frameworks}
\begin{itemize}
\item Backend runtime: Node.js 18+, TypeScript, Express
\item Data layer: PostgreSQL (Supabase), OpenSearch (BM25 + kNN), Redis + BullMQ
\item Storage layer: Cloudflare R2 and Supabase storage
\item Intelligence layer: Fashion-CLIP (ONNX), YOLO, BLIP, Gemini-assisted intent/attribute extraction, XGBoost ranker, Vertex AI try-on
\item Frontend clients: marketplace and admin applications
\item Deployment options: monolith mode, split service roles, Docker Compose, Cloud Run
\end{itemize}

\section{4.3 Feature Implementation by Module}
\subsection{4.3.1 Discover and Semantic Search}
Main entry points:

\begin{itemize}
\item \texttt{GET /search}
\item \texttt{GET /products/search}
\end{itemize}

Pipeline behavior:

\begin{enumerate}
\item Parse query and filters.
\item Run NLP-oriented query processing (intent/entity/negation/conversational context where available).
\item Build hybrid OpenSearch query using lexical + vector signals.
\item Apply relevance and business-aware reranking.
\item Hydrate results with product metadata and return paginated payload.
\end{enumerate}

Engineering notes:

\begin{itemize}
\item Search is mounted under \texttt{/search} (not \texttt{/api/search}).
\item Faceted browsing is exposed through \texttt{/products/facets}.
\item Trending/autocomplete endpoints provide retrieval-assistance UX.
\end{itemize}

\subsection{4.3.2 Visual Search (Single Image)}
Main entry points:

\begin{itemize}
\item \texttt{POST /products/search/image} (preferred integration path)
\item \texttt{POST /search/image} (alternate mount)
\end{itemize}

Pipeline behavior:

\begin{enumerate}
\item Accept multipart image input.
\item Generate image embeddings and optional attribute vectors.
\item Optionally extract caption/semantic hints (BLIP path when available).
\item Run vector retrieval on OpenSearch embedding fields.
\item Apply attribute-aware reranking (e.g., color/style alignment).
\item Return scored results with diagnostics metadata.
\end{enumerate}

Degradation behavior:

\begin{itemize}
\item If optional captioning/auxiliary signals are unavailable, pipeline continues with image-only retrieval.
\item Retrieval remains functional as long as core embedding path is available.
\end{itemize}

\subsection{4.3.3 Multi-Image and Multi-Vector Search}
Main entry points:

\begin{itemize}
\item \texttt{POST /search/multi-image}
\item \texttt{POST /search/multi-vector}
\end{itemize}

Multi-image behavior:

\begin{itemize}
\item Accepts up to multiple images with an optional natural-language prompt.
\item Uses prompt parsing and per-attribute embedding extraction.
\item Blends attribute-level signals into a weighted ranking decision.
\end{itemize}

Multi-vector behavior:

\begin{itemize}
\item Enables explicit weight control over dimensions (e.g., color/style/texture/material/pattern).
\item Useful for advanced users and experimental ranking control.
\end{itemize}

\subsection{4.3.4 Shop-the-Look (Detection + Per-Item Search)}
Main entry points:

\begin{itemize}
\item \texttt{POST /api/images/search}
\item \texttt{POST /api/images/search/url}
\item \texttt{POST /api/images/detect}
\end{itemize}

Extended endpoints in the same subsystem:

\begin{itemize}
\item \texttt{POST /api/images/analyze}
\item \texttt{POST /api/images/search/selective}
\item \texttt{POST /api/images/detect/url}
\item \texttt{POST /api/images/detect/batch}
\item status/metadata endpoints used for operational inspection
\end{itemize}

Purpose:

\begin{itemize}
\item Convert one outfit image into multiple actionable shopping paths.
\item Preserve item-level granularity so users can shop each detected garment separately.
\end{itemize}

Input contract patterns:

\begin{itemize}
\item multipart upload mode: direct image file
\item URL mode: remote image reference
\item selective mode: user-guided item/region emphasis
\item optional context: session/user identifiers, filter inheritance, ranking toggles
\end{itemize}

Detailed pipeline:

\begin{enumerate}
\item Request parsing and validation:
\end{enumerate}
\begin{itemize}
\item validate image presence and size/type constraints
\item normalize request-level options (limits, filters, mode flags)
\end{itemize}
\begin{enumerate}
\item Detection stage:
\end{enumerate}
\begin{itemize}
\item run detector and produce bounding boxes + labels + confidences
\item optionally return detect-only response for UI overlays/debug
\end{itemize}
\begin{enumerate}
\item Region preparation:
\end{enumerate}
\begin{itemize}
\item crop or isolate regions of interest per detection
\item generate region-level descriptors for retrieval
\end{itemize}
\begin{enumerate}
\item Per-detection retrieval:
\end{enumerate}
\begin{itemize}
\item run similarity search independently per detected item
\item enforce retrieval caps per detection to bound response size
\end{itemize}
\begin{enumerate}
\item Per-detection reranking:
\end{enumerate}
\begin{itemize}
\item combine embedding proximity with optional attribute/metadata relevance
\item apply relevance floors and remove obviously weak candidates
\end{itemize}
\begin{enumerate}
\item Diversity and deduplication:
\end{enumerate}
\begin{itemize}
\item collapse highly similar/variant-near results when enabled
\item avoid repetitive outputs in the same detection group
\end{itemize}
\begin{enumerate}
\item Grouped response assembly:
\end{enumerate}
\begin{itemize}
\item return one result bucket per detection (label, confidence, products)
\item include processing metadata and ranking diagnostics
\end{itemize}

Session and personalization behavior:

\begin{itemize}
\item Session context can propagate into item-level retrieval to preserve user intent continuity.
\item User context can enable personalized boosting when profile signals are available.
\item Fail-safe policy: missing session/profile data does not break base retrieval.
\end{itemize}

Quality controls:

\begin{itemize}
\item confidence thresholds for detection acceptance
\item bounded per-detection result counts
\item optional category/audience filters to reduce contamination
\item resilience paths for sparse retrieval (fallback query behavior)
\end{itemize}

Failure and degradation model:

\begin{itemize}
\item invalid input -> validation error response
\item detection service unavailable -> endpoint returns controlled failure (or detect-disabled path)
\item partial-stage failure in one detection bucket should not invalidate all successful buckets when recoverable
\end{itemize}

Output behavior:

\begin{itemize}
\item response returns grouped product suggestions by detected fashion item
\item each candidate can include score/reason-style metadata when available
\item diagnostics may include detection counts, dropped candidates, collapse metrics, and ranking mode flags
\end{itemize}

\subsection{4.3.5 Compare and Decision Intelligence}
Main entry points:

\begin{itemize}
\item \texttt{POST /api/compare} (primary decision-intelligence path)
\item \texttt{POST /api/compare/decision} (adapter path)
\item supporting quality/price/review/enhanced endpoints
\end{itemize}

Important additional compare endpoints:

\begin{itemize}
\item \texttt{GET /api/compare/quality/:productId}
\item \texttt{POST /api/compare/analyze-text}
\item \texttt{GET /api/compare/price/:productId}
\item \texttt{GET /api/compare/baseline/:category}
\item \texttt{POST /api/compare/admin/compute-baselines}
\item \texttt{GET /api/compare/tooltips}
\item review and enhanced compare endpoints for sentiment/logistics enrichment
\end{itemize}

Purpose:

\begin{itemize}
\item Move beyond "cheapest vs expensive" by giving structured purchase decisions.
\item Support direct alternatives, scenario-based choices, and outfit-context comparisons.
\end{itemize}

\texttt{POST /api/compare} full pipeline:

\begin{enumerate}
\item Input normalization:
\end{enumerate}
\begin{itemize}
\item accept IDs from JSON body or form-data
\item coerce arrays/strings into normalized unique numeric IDs
\end{itemize}
\begin{enumerate}
\item Request schema validation:
\end{enumerate}
\begin{itemize}
\item enforce 2..5 products
\item normalize aliases for goal, mode, occasion, and auxiliary signals
\end{itemize}
\begin{enumerate}
\item Product loading and profile shaping:
\end{enumerate}
\begin{itemize}
\item fetch product rows and normalize attributes
\item derive comparison-ready signals (price/text/style/metadata proxies)
\end{itemize}
\begin{enumerate}
\item Comparability gate:
\end{enumerate}
\begin{itemize}
\item block invalid cross-domain sets
\item enforce audience/category coherence rules where required
\end{itemize}
\begin{enumerate}
\item Mode resolution:
\end{enumerate}
\begin{itemize}
\item \texttt{direct\_head\_to\_head} for same-type alternatives
\item \texttt{scenario\_compare} for subtype/context differences
\item \texttt{outfit\_compare} for cross-slot set logic
\end{itemize}
\begin{enumerate}
\item Multi-axis scoring:
\end{enumerate}
\begin{itemize}
\item compute normalized dimensions: value, quality, style, risk, occasion
\item include practical vs expressive decision signals
\end{itemize}
\begin{enumerate}
\item Goal-aware weighting:
\end{enumerate}
\begin{itemize}
\item rebalance dimension weights by objective (\texttt{best\_value}, \texttt{premium\_quality}, \texttt{style\_match}, etc.)
\item apply occasion mismatch penalties when user context requires it
\end{itemize}
\begin{enumerate}
\item Winner and confidence decision:
\end{enumerate}
\begin{itemize}
\item determine leaders per context and overall
\item classify confidence (\texttt{clear\_choice}, \texttt{leaning\_choice}, \texttt{toss\_up})
\item optionally trigger "why-not-both" when margin is small and complementarity is high
\end{itemize}
\begin{enumerate}
\item Response serialization:
\end{enumerate}
\begin{itemize}
\item map engine output to stable frontend decision contract
\item include reasoning blocks, contextual winners, and confidence narrative
\end{itemize}

Legacy/adapter path (\texttt{POST /api/compare/decision}):

\begin{itemize}
\item exists for compatibility with older compare contract clients
\item runs legacy-style decision flow and adapts output to newer schema
\item useful during migration but not the primary strategic path
\end{itemize}

Heuristic and analytical sub-features:

\begin{itemize}
\item text quality analysis (description/detail richness)
\item price risk/anomaly analysis
\item image originality/quality proxies
\item return/shipping/policy interpretation
\item optional enhanced logistics/merchant/review intelligence
\end{itemize}

Observability and diagnostics:

\begin{itemize}
\item decision pipeline emits stage-level operational events
\item useful KPIs include mode distribution, fallback usage rate, toss-up frequency, and data quality warnings
\end{itemize}

Error model:

\begin{itemize}
\item invalid request -> 400
\item products not found -> 404
\item insufficient product data -> 422
\item unexpected engine failure -> 500
\end{itemize}

Response quality safeguards:

\begin{itemize}
\item comparability checks prevent misleading comparisons
\item confidence labeling prevents overconfident recommendations on ambiguous sets
\item mode-aware narration avoids forcing a single-winner message in outfit-context scenarios
\end{itemize}

\subsection{4.3.6 CompleteStyle for Products (This is NOT Complete-Look)}
Main entry points:

\begin{itemize}
\item \texttt{GET /products/:id/complete-style}
\item \texttt{POST /products/complete-style}
\item \texttt{POST /products/complete-style/try-on}
\end{itemize}

Important distinction:

\begin{itemize}
\item \texttt{complete-style} = complementary product recommendation from a source anchor product.
\item It is a recommendation pipeline focused on completing a product-centered outfit suggestion.
\item It is explicitly documented as a separate pipeline from \texttt{complete-look}.
\end{itemize}

Unified CompleteStyle pipeline stages:

\begin{enumerate}
\item Request normalization:
\end{enumerate}
\begin{itemize}
\item options merging and clamping (\texttt{maxPerCategory}, \texttt{maxTotal}, price filters, brand exclusion/preference, audience hints)
\end{itemize}
\begin{enumerate}
\item Source anchor resolution:
\end{enumerate}
\begin{itemize}
\item resolve catalog source product or payload product
\end{itemize}
\begin{enumerate}
\item Source understanding:
\end{enumerate}
\begin{itemize}
\item detect category, infer style profile, infer audience metadata
\end{itemize}
\begin{enumerate}
\item Candidate retrieval:
\end{enumerate}
\begin{itemize}
\item expand complementary categories
\item retrieve oversampled candidates with hard filters + soft boosts
\end{itemize}
\begin{enumerate}
\item Candidate scoring:
\end{enumerate}
\begin{itemize}
\item hard gates (reject invalid/incompatible/self rows)
\item soft compatibility signals (style, color harmony, occasion, attributes, brand preference)
\end{itemize}
\begin{enumerate}
\item Reranking and confidence filtering:
\end{enumerate}
\begin{itemize}
\item ML ranker when available, heuristic fallback otherwise
\end{itemize}
\begin{enumerate}
\item Wardrobe fusion (user-aware):
\end{enumerate}
\begin{itemize}
\item mark owned items
\item inject owned extras in sparse categories
\item preserve relevance-first ordering
\end{itemize}
\begin{enumerate}
\item Coverage balancing:
\end{enumerate}
\begin{itemize}
\item prioritize essentials, enforce category caps, dedupe, controlled relaxation
\end{itemize}
\begin{enumerate}
\item Response mapping:
\end{enumerate}
\begin{itemize}
\item return grouped recommendations with reasons/scores and completion metadata
\end{itemize}

Resilience model:

\begin{itemize}
\item Ranker unavailable -> heuristic fallback.
\item Visual signal unavailable -> disable visual contribution only.
\item Wardrobe unavailable -> continue without wardrobe merge.
\end{itemize}

\subsection{4.3.7 Complete-Look for Wardrobe and Catalog Anchors}
Main entry point:

\begin{itemize}
\item \texttt{POST /api/wardrobe/complete-look}
\end{itemize}

Supported request modes:

\begin{itemize}
\item \texttt{item\_ids} (wardrobe anchors)
\item \texttt{product\_ids} (catalog anchors)
\end{itemize}

Core purpose:

\begin{itemize}
\item Build coherent multi-item outfit sets, not just per-item complements.
\item Output includes \texttt{suggestions}, \texttt{outfitSets}, and \texttt{missingCategories}.
\end{itemize}

Complete-Look pipeline stages:

\begin{enumerate}
\item Anchor intake and mode detection (\texttt{wardrobe} vs \texttt{catalog-product})
\item Slot inference (tops, bottoms, shoes, dresses, outerwear, bags, accessories)
\item Gap detection (essential + complement slots; weather-aware constraints)
\item Candidate retrieval per missing slot:
\end{enumerate}
\begin{itemize}
\item strict filters first
\item vector/lexical retrieval blend
\item controlled recall relaxation if sparse
\end{itemize}
\begin{enumerate}
\item Candidate scoring with stylist features:
\end{enumerate}
\begin{itemize}
\item category compatibility
\item color harmony
\item style alignment
\item pattern/material/formality alignment
\end{itemize}
\begin{enumerate}
\item Fashion-aware reranking:
\end{enumerate}
\begin{itemize}
\item blend of fashion score and base retrieval score
\end{itemize}
\begin{enumerate}
\item Outfit set builder:
\end{enumerate}
\begin{itemize}
\item enumerate cross-slot combinations from top candidates
\item score pairwise coherence between recommended pieces
\end{itemize}
\begin{enumerate}
\item Coherence gate:
\end{enumerate}
\begin{itemize}
\item remove low-coherence set combinations
\end{itemize}
\begin{enumerate}
\item Return top outfit sets with explanatory signals.
\end{enumerate}

Why Complete-Look is distinct:

\begin{itemize}
\item CompleteStyle is a product-centric recommendation pipeline.
\item Complete-Look is a set-construction stylist pipeline with pairwise outfit coherence gating.
\item Complete-Look includes explicit slot and set-level reasoning that cannot be reduced to simple nearest-neighbor retrieval.
\end{itemize}

\subsection{4.3.8 Recommendations, Wardrobe, and Styling Intelligence}
Recommendations:

\begin{itemize}
\item \texttt{GET /products/:id/recommendations}
\item \texttt{POST /products/recommendations/batch}
\end{itemize}

Wardrobe core:

\begin{itemize}
\item CRUD item management and profile computation
\item gap detection and compatibility graph endpoints
\item outfit suggestions and coherence scoring
\item layering analysis and weather checks
\item photo analysis and re-analysis pipelines
\end{itemize}

Wardrobe intelligence value:

\begin{itemize}
\item Bridges marketplace recommendations with user-owned inventory.
\item Enables style-consistent, practicality-aware suggestions.
\end{itemize}

Detailed recommendation flow:

\begin{enumerate}
\item Candidate generation:
\end{enumerate}
\begin{itemize}
\item retrieve similarity candidates from vector/text signals
\item include category-compatible products
\end{itemize}
\begin{enumerate}
\item Feature engineering:
\end{enumerate}
\begin{itemize}
\item style/color/price similarity features
\item optional image-derived and historical behavior signals
\end{itemize}
\begin{enumerate}
\item Ranking stage:
\end{enumerate}
\begin{itemize}
\item ML ranker path (XGBoost service) when available
\item heuristic fallback weights when ranker unavailable
\end{itemize}
\begin{enumerate}
\item Post-processing:
\end{enumerate}
\begin{itemize}
\item deduplicate near-identical results
\item enforce caps/diversity and optional policy filters
\end{itemize}
\begin{enumerate}
\item Response:
\end{enumerate}
\begin{itemize}
\item ranked items plus optional score/reason metadata
\end{itemize}

Wardrobe detailed capabilities:

\begin{itemize}
\item profile computation from owned items (style/color/category distributions)
\item compatibility graph creation for pairwise item fit
\item wardrobe gap analysis for missing essentials
\item outfit suggestion generation from owned inventory
\item image/photo analysis for item auto-categorization
\item coherence scoring for user-defined outfits
\item layering analysis with weather-awareness constraints
\end{itemize}

\subsection{4.3.9 Authentication and User Management}
Main entry points:

\begin{itemize}
\item \texttt{POST /api/auth/signup}
\item \texttt{POST /api/auth/login}
\item \texttt{POST /api/auth/refresh}
\item \texttt{POST /api/auth/logout}
\item \texttt{GET/PATCH /api/auth/me}
\end{itemize}

Auth flow:

\begin{enumerate}
\item Signup/login validates credentials and account status.
\item Access + refresh tokens are issued with bounded lifetimes.
\item Protected routes rely on auth middleware.
\item Refresh endpoint rotates token sessions.
\item Logout blacklists refresh token to prevent reuse.
\end{enumerate}

Current limitations:

\begin{itemize}
\item no email verification
\item no forgot/reset password flow
\item strong JWT secret must be enforced by production environment configuration
\end{itemize}

\subsection{4.3.10 Cart and Favorites}
Cart entry points:

\begin{itemize}
\item \texttt{GET /api/cart}
\item \texttt{POST /api/cart}
\item \texttt{PATCH /api/cart/:productId}
\item \texttt{DELETE /api/cart/:productId}
\item \texttt{DELETE /api/cart/clear}
\end{itemize}

Cart behavior:

\begin{enumerate}
\item Validate authenticated user context.
\item Upsert/add item with quantity semantics.
\item Join product data for display-ready cart payload.
\item Recompute totals server-side for consistency.
\end{enumerate}

Favorites entry points:

\begin{itemize}
\item \texttt{GET /api/favorites}
\item \texttt{POST /api/favorites/toggle}
\item \texttt{GET /api/favorites/check/:productId}
\item \texttt{POST /api/favorites/check} (batch)
\end{itemize}

Favorites behavior:

\begin{itemize}
\item toggle persistence by user/product
\item supports single and batch existence checks for fast UI rendering
\end{itemize}

\subsection{4.3.11 Ingest Pipeline and Labeling Pipeline}
Ingest entry points:

\begin{itemize}
\item \texttt{POST /api/ingest/image}
\item \texttt{GET /api/ingest/:jobId}
\end{itemize}

Ingest processing flow (worker-based):

\begin{enumerate}
\item Accept upload and register ingest job.
\item Queue job through Redis/BullMQ.
\item Worker downloads/validates image payload.
\item Compute embeddings and run detection.
\item Produce derivative artifacts and metadata.
\item Persist generated records and mark job status.
\end{enumerate}

Design characteristics:

\begin{itemize}
\item asynchronous execution to protect API latency
\item status polling contract for frontend reliability
\item fault isolation between ingestion and user-facing request path
\end{itemize}

Labeling entry points:

\begin{itemize}
\item task retrieval/assignment/submission/skip/stats endpoints under \texttt{/api/labeling}
\end{itemize}

Labeling purpose:

\begin{itemize}
\item provide active-learning feedback loops for recommendation quality
\item collect supervised signals for model refinement
\end{itemize}

\subsection{4.3.12 Image Analysis Subsystem}
Main entry points:

\begin{itemize}
\item \texttt{POST /api/images/analyze}
\item \texttt{POST /api/images/search}
\item \texttt{POST /api/images/search/selective}
\item \texttt{POST /api/images/search/url}
\item \texttt{POST /api/images/detect}
\item \texttt{POST /api/images/detect/url}
\item \texttt{POST /api/images/detect/batch}
\end{itemize}

Subsystem behavior:

\begin{enumerate}
\item detection-only mode for diagnostics and UI overlays
\item search mode for retrieval over detected items
\item selective mode for user-directed region/item intent
\item URL-based mode for remote image workflows
\item batch mode for throughput-oriented image processing
\end{enumerate}

\subsection{4.3.13 Admin and Operations Endpoints}
Admin scope:

\begin{itemize}
\item product moderation (hide/unhide/flag/unflag)
\item canonical grouping and catalog hygiene
\item operational job triggers and stats endpoints
\end{itemize}

Security posture:

\begin{itemize}
\item admin endpoints are protected by auth + admin-role middleware
\item operations remain auditable through structured logs and metrics
\end{itemize}

\subsection{4.3.14 Metrics, Health, and Observability}
Health endpoints:

\begin{itemize}
\item \texttt{GET /health/live}
\item \texttt{GET /health/ready}
\end{itemize}

Metrics endpoint:

\begin{itemize}
\item \texttt{GET /metrics} (Prometheus format)
\end{itemize}

Observability goals:

\begin{itemize}
\item enable readiness/liveness orchestration
\item expose request volume and latency instrumentation
\item support incident triage and deployment confidence checks
\end{itemize}

\subsection{4.3.15 Data and Search Infrastructure Internals}
PostgreSQL responsibilities:

\begin{itemize}
\item source-of-truth for users, catalog metadata, wardrobe, try-on jobs, and workflow tables
\item migration-driven schema evolution
\end{itemize}

OpenSearch responsibilities:

\begin{itemize}
\item vector and lexical retrieval index
\item hybrid query execution and ranking candidate retrieval
\end{itemize}

Redis/BullMQ responsibilities:

\begin{itemize}
\item queueing and background processing
\item decoupling heavy compute from synchronous API paths
\end{itemize}

Storage responsibilities:

\begin{itemize}
\item image and generated asset persistence
\item stable URL/object references for downstream services
\end{itemize}

\subsection{4.3.16 Virtual Try-On}
Main entry points under \texttt{/api/tryon}:

\begin{itemize}
\item submit (\texttt{/}, \texttt{/from-wardrobe}, \texttt{/from-product}, \texttt{/batch})
\item status/history (\texttt{/:id}, \texttt{/history}, \texttt{/saved})
\item lifecycle management (\texttt{cancel}, \texttt{delete}, \texttt{save})
\end{itemize}

Operational model:

\begin{itemize}
\item async job lifecycle (\texttt{pending -> processing -> completed/failed})
\item immediate acceptance on submit and polling-based retrieval
\item dependency on Vertex AI credentials/project configuration
\end{itemize}

Detailed try-on flow:

\begin{enumerate}
\item Accept person/garment sources (file or product/wardrobe references).
\item Validate user context and rate limit boundaries.
\item Create async job record and return accepted response.
\item Background processor calls external try-on model endpoint.
\item Persist generated result artifacts and status transitions.
\item Expose history/saved management for user experience continuity.
\end{enumerate}

Failure handling:

\begin{itemize}
\item configuration and migration errors are surfaced with actionable error codes
\item job-level failures do not crash API process; they transition to failed state
\end{itemize}

\section{4.4 Security, Reliability, and Operations Controls}
Security controls:

\begin{itemize}
\item Helmet and CORS middleware
\item API rate limiting and bounded payload parsing
\item JWT access/refresh authentication
\item refresh token blacklist logout behavior
\item route-level protection for user and admin scopes
\end{itemize}

Reliability controls:

\begin{itemize}
\item health probes (\texttt{/health/live}, \texttt{/health/ready})
\item Prometheus metrics (\texttt{/metrics})
\item graceful startup/degradation around model/index initialization
\item queue-based asynchronous processing for heavy jobs
\end{itemize}

Operational controls:

\begin{itemize}
\item deployment runbooks and env contracts
\item migration-driven schema evolution
\item index management scripts and reindex support
\item role-based service deployment options (monolith or split-role modes)
\end{itemize}

\chapter{5. Testing and Evaluation}
\section{5.1 Testing Strategy}
The system applies mixed validation layers:

\begin{enumerate}
\item Unit tests for relevance, compatibility, and feature-engine logic.
\item Service-level tests for decision/ranking behavior.
\item Scripted evaluations for search quality and image relevance.
\item Type-check and regression matrix validation for stylist pipelines.
\end{enumerate}

Representative validation evidence includes:

\begin{itemize}
\item search relevance utility tests
\item decision-intelligence feature tests
\item complete-look matrix validation
\item TypeScript compile-time integrity checks
\end{itemize}

\section{5.2 Evaluation Methodology}
Evaluation is performed along practical engineering dimensions:

\begin{itemize}
\item Retrieval quality for text, image, and multi-image queries
\item Robustness of compare decisions under heterogeneous product sets
\item Outfit completion quality at both per-item and set level
\item Stability under optional component failure
\item Operational readiness (health/metrics/deployability)
\end{itemize}

\section{5.3 Feature-Specific Evaluation Notes}
\subsection{5.3.1 Search and Discovery}
\begin{itemize}
\item Hybrid retrieval improves recall over lexical-only approaches.
\item Attribute-aware reranking improves semantic alignment in visual search.
\item Session/personalization propagation improves contextual relevance continuity.
\end{itemize}

\subsection{5.3.2 Compare Decision Intelligence}
\begin{itemize}
\item Multi-axis scoring improves explainability versus single-score ranking.
\item Confidence classification (\texttt{clear\_choice}, \texttt{leaning\_choice}, \texttt{toss\_up}) enables safer UX messaging.
\item Comparability gates reduce misleading cross-domain comparisons.
\end{itemize}

\subsection{5.3.3 CompleteStyle (Product)}
\begin{itemize}
\item Category expansion + scoring + balancing improves recommendation completeness.
\item Wardrobe-aware fusion provides practical personalization.
\item Controlled relaxation preserves non-empty useful outputs without violating hard safety gates.
\end{itemize}

\subsection{5.3.4 Complete-Look (Wardrobe/Catalog)}
\begin{itemize}
\item Slot inference and missing-category logic improve outfit assembly quality.
\item Pairwise coherence scoring reduces incoherent combinations that would pass item-level relevance.
\item Coherence gating enforces stylist consistency at set-level output.
\end{itemize}

\section{5.4 Known Gaps and Evaluation Constraints}
\begin{itemize}
\item Checkout/payment/order lifecycle is not fully implemented.
\item Email verification and password reset are absent.
\item Test distribution is uneven across all modules.
\item Some environment/migration setup errors can degrade specific features if not operationally enforced.
\item Privacy policy artifacts are less formal than technical controls.
\end{itemize}

\chapter{6. Entrepreneurial / Innovation Aspects}
\section{6.1 Market Relevance}
The platform addresses a clear local need: unifying fragmented fashion inventory and reducing search friction in a multi-store ecosystem.


\section{6.2 Innovation}
Innovation lies in end-to-end integration:

\begin{itemize}
\item Multimodal search and shop-the-look
\item Decision intelligence (compare + rank + complete style)
\item Wardrobe-aware recommendations
\item Cloud try-on and operationalized ML components
\end{itemize}

\section{6.3 Commercialization and Scalability}
The architecture supports expansion via:

\begin{itemize}
\item Additional vendors and catalogs
\item Region-level rollout
\item Service-role split deployment (API vs ML)
\item Feature-level monetization opportunities (premium try-on/styling intelligence)
\end{itemize}

\section{6.4 Ethical and Societal Considerations}
Ethical priorities include fair product representation, transparent recommendation logic, and responsible handling of user-uploaded images and identity-related data in try-on scenarios.


\chapter{7. Project Management and Teamwork}
\section{7.1 Timeline and Milestones}
Typical milestones:

\begin{enumerate}
\item Requirements and architecture definition
\item Core API and search foundation
\item Intelligence modules (recommendation/compare/wardrobe/try-on)
\item Testing, hardening, and documentation consolidation
\item Deployment-readiness and operational monitoring
\end{enumerate}

\section{7.2 Resource Management}
Key resources included cloud services, ML model runtime dependencies, search infrastructure, and structured documentation for maintainability.


\section{7.3 Team Roles}
Representative role split:

\begin{itemize}
\item Backend/AI integration engineering
\item Frontend/dashboard integration
\item Data/ML pipeline support
\item Quality assurance and release stabilization
\end{itemize}

\chapter{8. Results and Discussion}
\section{8.1 Objective Achievement}
The project successfully delivers a technically rich unified fashion discovery platform with multimodal retrieval, comparison, recommendation, wardrobe intelligence, and virtual try-on capabilities.


\section{8.2 Strengths}
\begin{itemize}
\item Broad feature integration under one API ecosystem
\item Practical use of modern AI/ML components
\item Scalable and modular architecture with clear separation of concerns
\item Mature operational signals (health, metrics, deployment docs)
\end{itemize}

\section{8.3 Limitations}
\begin{itemize}
\item Commerce completion capabilities remain partial (no full order lifecycle)
\item Some security/config defaults require strict production governance
\item Reliance on external ML and cloud services introduces operational coupling
\item Coverage and policy documentation can be further strengthened
\end{itemize}

\section{8.4 Contribution to Practice}
This work demonstrates how computer science concepts (information retrieval, machine learning, distributed systems, API engineering, and software architecture) can be integrated into a market-relevant product with real deployment pathways.


\chapter{9. Conclusion and Future Work}
\section{9.1 Conclusion}
This capstone implements a professional-grade fashion discovery platform that bridges the gap between inspiration and purchase intent. The system advances beyond simple search by combining multimodal retrieval, decision intelligence, and style-aware workflows in a cohesive architecture.


\section{9.2 Future Work}
Recommended next steps:

\begin{itemize}
\item Implement checkout, orders, and payment lifecycle
\item Add email verification and password recovery
\item Expand automated integration and end-to-end test suites
\item Introduce online experimentation (A/B testing) for ranking and UX decisions
\item Formalize privacy/data-retention policy and governance artifacts
\item Strengthen model lifecycle management (monitoring, retraining, drift control)
\end{itemize}

\chapter{References}
[1] Project root documentation and system overview: \texttt{README.md}

[2] Feature mapping and endpoint integration: \texttt{docs/FEATURES.md}

[3] API contract details: \texttt{docs/api-reference.md}

[4] Architecture and module guidelines: \texttt{docs/architecture.md}

[5] Implementation audit and known gaps: \texttt{docs/IMPLEMENTATION\_STATUS.md}

[6] Embeddings and retrieval pipelines: \texttt{docs/embeddings-and-search-pipelines.md}

[7] ML stack details: \texttt{docs/ml-models.md}

[8] Data model and migrations: \texttt{docs/database.md}

[9] Deployment guides: \texttt{docs/deployment.md}, \texttt{docs/deploy-cloud-run.md}


\chapter{Appendices}
\section{Appendix A: API Endpoint Domains}
\begin{itemize}
\item \texttt{/search}, \texttt{/products}, \texttt{/api/images}, \texttt{/api/compare}, \texttt{/api/wardrobe}, \texttt{/api/tryon}, \texttt{/api/auth}, \texttt{/api/cart}, \texttt{/api/favorites}, \texttt{/admin}, \texttt{/metrics}, \texttt{/health}
\end{itemize}

\section{Appendix B: Suggested Figures/Tables for Final DOCX}
\begin{itemize}
\item System architecture diagram (layered + module map)
\item Search pipeline flow (text/image/multi-image)
\item Database entity overview
\item Try-on async state machine
\item Feature-to-endpoint traceability matrix
\end{itemize}
\end{document}
